% Reused from chapters/05-ospf.tex
\begin{labvisual}{Recall the chapter model}
\begin{labgraphic}
\begin{tikzpicture}[node distance=1.0cm and 1.5cm,
  db/.style={rectangle,rounded corners=1mm,draw=UTPTeal,fill=UTPPaleTeal,
             minimum width=2.5cm,minimum height=0.95cm,align=center,
             font=\scriptsize\color{UTPNavy}},
  ml/.style={font=\scriptsize\color{UTPMuted},align=center}]

\node[router] (r1) {R1};
\node[router,right=2.2cm of r1] (r2) {R2};
\node[router,right=2.2cm of r2] (r3) {R3};
\draw[link] (r1) -- (r2);
\draw[link] (r2) -- (r3);
\draw[link] (r1) to[bend right=32] (r3);

\node[db,below=1.15cm of r1] (d1) {identical\\ link-state DB};
\node[db,below=1.15cm of r2] (d2) {identical\\ link-state DB};
\node[db,below=1.15cm of r3] (d3) {identical\\ link-state DB};
\draw[<->,thick,draw=UTPTeal] (d1) -- node[ml,above]{flood} (d2);
\draw[<->,thick,draw=UTPTeal] (d2) -- node[ml,above]{flood} (d3);

\node[ml,below=0.45cm of d2,text width=9.5cm]
  {Same database everywhere. Each router then runs Dijkstra \textbf{from itself},
   so the routing tables all differ --- and all agree.};
\end{tikzpicture}
\end{labgraphic}\end{labvisual}
